Lors d'un tremblement de terre, deux types d'ondes se propagent dans la croûte terrestre : les ondes P (longitudinales) à \qty{\v{vP}}{km/s} et les ondes S (transversales) à \qty{\v{vS}}{km/s}. Une station sismique enregistre l'arrivée des ondes S \qty{\v{dt}}{s} après celle des ondes P.

\begin{enumerate}[label=\alph*)]
    \item À quelle distance de la station le séisme s'est-il produit ?
    \item Combien de temps après le séisme les ondes P ont-elles atteint la station ?
    \item Pourquoi les ondes S ne peuvent-elles pas traverser le noyau externe liquide de la Terre, alors que les ondes P le peuvent ?
\end{enumerate}
